\documentclass[12pt]{article}
\usepackage[utf8]{inputenc}
\usepackage[T1]{fontenc}
\usepackage{hyperref}
\title{Malay, Central Language Resources for K-12 Curriculum Development}
\author{Compiled by Claude Sonnet 5.5 (claude-sonnet-5-5) for Andr\'as Kornai}
\date{}
\begin{document}
\maketitle

\section{Introduction}
ISO 639-3 pse (Central Malay) corresponds to South Barisan Malay, also called Middle Malay: a group of closely related Malayic varieties in southwestern Sumatra (Bengkulu, South Sumatra, Lampung), with about 1.6 million speakers according to Wikipedia. Varieties include Bengkulu Malay, Serawai, Besemah (Pasemah), Ogan, Enim, Lematang, Semende and others. This language must not be confused with Ambonese Malay (abs). Search results turned up very little educational material specific to it.

\section{Primary Educational Resources}
No textbooks, primers or school programmes for this language were found in search results.

\section{Linguistic and Reference Materials}
\begin{itemize}
\item Wikipedia, South Barisan Malay: classification, dialect list and speaker estimate. \url{https://en.wikipedia.org/wiki/South_Barisan_Malay}
\item Wikipedia ISO 639:pse page for the same entry. \url{https://wikipedia.com/wiki/ISO_639:pse}
\end{itemize}

\section{Religious and Community Resources}
\begin{itemize}
\item Joshua Project profile of the Serawai, a Malay group in South Bengkulu and Seluma. \url{https://www.joshuaproject.net/people_groups/14863}
\item Joshua Project profile of the Bengkulu. \url{https://joshuaproject.net/people_groups/10791}
\item Joshua Project language page for pse. \url{https://www.joshuaproject.net/index.php/languages/pse}
\end{itemize}

\section{Digital and Audio Resources}
\begin{itemize}
\item Global Recordings Network, Malay, Central: Bengkulu language page (availability of recordings not verified). \url{https://www.globalrecordings.net/en/language/25934}
\end{itemize}

\section{Curriculum Development Context}
The language is a cluster of regional varieties, taught as a subject in some schools according to the reference description, but no curricula were found. Indonesian is the school language; local-content (muatan lokal) materials would have to be developed or located through regional education offices, which these searches did not reach.
\end{document}